\section{Position Representation：Self-Attention 为什么必须补回顺序}

Self-attention 对输入 permutation equivariant：若同时打乱 token 和位置，纯 attention 无法知道原顺序。语言、音乐和图像都依赖相对位置，因此模型必须注入 position signal。原始 Transformer 的 learned/sinusoidal encoding 是起点，后续 relative embedding、relative bias 和 rotation 直接把距离写入 interaction。

\subsection{Absolute Position：Learned 与 Sinusoidal}

原始方法把 position vector $p_i$ 加到 token embedding $x_i$。Sinusoidal encoding 的一组维度为：
\[
p_{i,2k}=\sin\!\left(i/10000^{2k/d}\right),\qquad
p_{i,2k+1}=\cos\!\left(i/10000^{2k/d}\right).
\]
其中 $i$ 是位置，$k$ 是频率索引，$d$ 是 model width。它无需学习、可计算任意位置，但 model 是否真正外推到训练长度之外仍取决于 attention behavior 与训练分布。

\lecturefigure{slide-27-evolution-position-encodings.jpg}{Permutation invariance 与原始 learned/sinusoidal position encoding}{官方录像恢复幻灯片 V027；00:33:45}

\begin{knowledgebox}{读图：Position signal 应该回答“相对多远”，不只回答“绝对第几”}
右侧原始架构在底部加 position encoding。先看绝对 index 能否自然表达 $i-j$，再看训练长度之外是否有稳定结构；learned embedding 容易受 max length 限制，sinusoid 有解析形式却不保证模型学会可靠 extrapolation。
\end{knowledgebox}

\teachervoice{讲者说原始 sinusoidal/learned position 无法很好捕获 relative distance 与 extrapolation，这成为后续演化重点。调整 base frequency 可以改变长度行为，但 position scheme 不是一个独立魔法旋钮，仍要与训练 context 和 attention scale 配合。}

\subsection{Relative Embedding、Bias 与 Rotation}

Relative 方法把距离 $r=i-j$ 直接加入 score 或 value。一个简化 relative-key attention 为：
\[
s_{ij}=q_i^\top(k_j+a_{i-j}),
\]
其中 $a_{i-j}$ 是相对距离 embedding。Relative bias 直接向 $s_{ij}$ 加标量，rotary position embedding（RoPE）则按位置旋转 query/key，使内积包含相对 phase；三者对参数、外推和实现的取舍不同。

\lecturefigure{slide-28-position-encoding-variants.jpg}{Relative embedding、relative bias 与 relative rotation}{官方录像恢复幻灯片 V028；00:36:45}

\begin{knowledgebox}{术语消化：Position 方案的接口差异}
\begin{tabularx}{\textwidth}{@{}L{0.20\textwidth}X X@{}}
\toprule
方案 & 把位置放在哪里 & 主要取舍 \\
\midrule
Absolute embedding & 加到 token hidden state & 简单，但距离需间接学习 \\
Relative embedding/bias & 加入 pairwise attention score/value & 直接表达距离，需处理距离范围 \\
Rotary encoding & 旋转 query/key subspace & 内积自然含相对 phase，外推依赖频率设计 \\
\bottomrule
\end{tabularx}
\end{knowledgebox}

\subsection{Music Transformer：相对位置如何改善长结构}

音乐生成要求模型维持节奏、主题和重复结构，token 间依赖可能跨越很长时间。Music Transformer 使用 relative attention，使模型更容易按间隔与重复关系读取历史；它展示 position mechanism 可以改变 qualitative coherence，而不只是 benchmark 数字。

\lecturefigure{slide-29-music-transformer.jpg}{Music Transformer 用 relative attention 建模长程音乐结构}{官方录像恢复幻灯片 V029；00:37:12}

\begin{knowledgebox}{读图：先看 relation，再听/看 outcome}
图中强调 relative attention 的 score decomposition。先追踪当前 note 与历史 note 的时间间隔，再比较 vanilla attention；relative signal 提供结构先验，但最终音乐质量还依赖 tokenization、training corpus、sampling 和评价方法。
\end{knowledgebox}

\lecturefigure{slide-30-music-transformer-demo.jpg}{Music Transformer 与其他模型的 qualitative sample 页面}{官方录像恢复幻灯片 V030；00:37:18}

\begin{warningbox}{Demo 是证据，但证据强度有限}
网页样例可帮助听众感受 repetition、coherence 和 drift，却容易受 cherry-picking、播放片段和主观偏好影响。应把它与 held-out likelihood、结构指标、blind listening study 和多 seed 样例结合，而不是把一段好听音乐当作普遍胜利。
\end{warningbox}

\teachervoice{讲者播放/展示音乐 sample，是为了让 relative position 的影响可感知；他没有把 demo 当作完整 metric。课堂保留这种定性证据，同时必须明确它只能提出“长程结构更好”的观察，不足以确定唯一原因。}

\subsection{本章小结}

Self-attention 本身不知道顺序。Absolute encoding 让模型看到 index，relative/bias/rotation 让 pairwise interaction 直接感知距离，Music Transformer 则提供长结构应用。Position representation 是 attention interface 的核心组成，而非输入前的一次装饰性相加。

